\chapter{第一批细胞可能怎样出现}

\begin{kaichang}
去河边或海边捡两样东西：一颗被水磨圆的鹅卵石，和一只在石头缝里爬的小虫（潮虫、小螃蟹都行）。把它们放在手心里比一比：一个一动不动，放一百年还是那颗石头；另一个爬来爬去，会吃东西、会躲你，将来还会养出一窝小虫。

再想想它们其实是什么：第 7 章说过，万物都由周期表上那一百来种元素的原子组成。石头主要是硅和氧，虫子主要是碳、氢、氧、氮——用的都是同一盒积木。为什么同一盒积木搭出来的东西，一个永远“死”着，另一个却“活”了？

更刺激的问题在后面：40 亿年前的地球上只有石头、海洋和空气，连一只细菌都没有。那么，第一个“活”的东西是从哪儿冒出来的？先猜一个你自己的答案。这一章我会把科学家手里的证据和空白都摊给你看——你会读到很多“可能”和“还不知道”。先说清楚，这不是卖关子，而是这个问题最真实的现状。
\end{kaichang}

\section{现象层：岩石记账本上的时间}

先盘点那些真的摆在那儿、可以测量的东西。

第一件：地球的年龄。用第 8 章讲过的放射性同位素测年，最古老的陨石和月岩都把答案指向同一个数字：太阳系和地球大约形成于 45.4 亿年前。这不是猜的，是好几种独立的“时钟”（比如铀变成铅）互相对出来的。

第二件：水来得比想象早。澳大利亚西部杰克山出土的锆石小晶体，有的年龄高达 44 亿年。它们内部的氧同位素比例暗示：这些晶体形成时，地表附近可能已经有液态水了。也就是说，地球诞生后没几亿年，海洋就可能已经出现。

第三件：生命登场“快”得反常。澳大利亚皮尔巴拉地区保存着约 35 亿年前的叠层石——一层一层、像卷心菜剖面一样的石头，被认为是远古微生物群落一遍遍生长、黏住泥沙留下的遗迹。格陵兰还有约 37 亿年前的疑似叠层石，以及 38 亿年前岩石里“像被生物加工过”的碳同位素信号，但这些更早的证据至今有争议，地质学家还在争。即便只认最保守的 35 亿年，结论也够惊人：从“海洋出现”到“生命留下痕迹”，中间只隔了几亿年，甚至更短。在地质尺度上，这几乎等于水一凑齐，生命很快就上了场。

第四件：更早的记录永远缺页。地球表面在板块运动中不断被回收、熔化、重铸，40 亿年前的海洋和地壳几乎荡然无存。所以“生命到底在哪一刻、在哪里、以什么方式开始”，最直接的物证早被地球自己销毁了。这正是本章注定充满“可能”二字的根本原因。

还有一个 1977 年才登场的重要现场：深海热泉。那一年，科学家乘潜水器在加拉帕戈斯附近的洋底裂谷，第一次看到“黑烟囱”——滚烫的、裹着矿物颗粒的黑水从海底喷口涌出。更令人目瞪口呆的是，喷口周围挤满管虫、蛤蜊和虾，整套生态系统不靠一缕阳光，全靠细菌“吃”硫化氢、氢气这类化学物质过活。原来生命不一定需要太阳，岩石和水的化学反应就能养活一座城市。这给“生命可能起步于海底”的猜想，提供了一个活着的样板。

\section{粒子层：先把零件备齐}

把镜头推到分子尺度，先列一张开工清单。从第 45 到 55 章你已经认识了生命的主要分子：蛋白质由氨基酸串成（第 48 章），膜主要由脂质搭成（第 47 章），遗传信息存在核苷酸链里，糖既当燃料又当材料（第 46 章）。问题因此变得非常具体：\textbf{在没有任何生物的地球上，氨基酸、脂质分子、核苷酸这些零件，能从无机环境里自己冒出来吗？}

今天我们知道，至少有三条来源，而且条条有实证。

第一条，天空。早期地球没有臭氧层，紫外线直射地表；火山喷发、闪电、陨石撞击也比今天频繁得多。在这种高能环境里，简单气体分子之间完全可能发生反应，生成更复杂的有机小分子。证据故事栏目里，马上会讲那场著名的实验。

第二条，海底。深海热泉口不断涌出热的、富含氢气和矿物颗粒的水流，遇上冰冷的海水和 \chem{CO_2}。热泉壁上疏松的矿物微孔像无数口并排的小反应锅，矿物表面本身又能当催化剂——回想第 30 章：催化剂不改变反应方向，只给反应铺一条好走的路。今天，科学家在热泉口确实检测到了简单有机物，也在实验室模拟热泉的条件下合成出了脂肪酸样的分子。

第三条，天外快递。1969 年落在澳大利亚的默奇森陨石，内部检出 70 多种氨基酸——其中大部分根本不是地球生物使用的那 20 种。这恰恰说明它们不是地球上的“污染”，而是太空里化学合成的产物。今天每年仍有成千上万吨星际尘埃落到地球，早期地球接收的“快递”只会更多。

三条来源指向同一个结论：\textbf{零件，不成问题。}有机小分子能在无生命的环境里形成，这一点已不再是猜想，而是测得到的事实。真正的难题在下一步。

\begin{zhengju}
1952 年，芝加哥大学的研究生斯坦利·米勒在导师哈罗德·尤里的支持下，搭了一套玻璃装置：一个烧瓶装水加热，当“海洋”；上面的大烧瓶装甲烷（\chem{CH_4}）、氨（\chem{NH_3}）、氢气（\chem{H_2}）和水蒸气，当“原始大气”；再插两根电极不停打出电火花，模拟闪电。整套装置密封循环，跑了一个星期。

一周后，“海洋”变成了红褐色。分析结果：里面生成了甘氨酸、丙氨酸等多种氨基酸——蛋白质的基本零件，居然从一锅无机气体里自己长了出来。1953 年论文发表，轰动世界：生命起源似乎从哲学问题，一夜之间变成了可以做实验的化学问题。

但这个故事的精彩处才刚开始。后来的研究指出：米勒假设的大气成分可能不对。越来越多的证据显示，早期地球大气以 \chem{CO_2} 和 \chem{N_2} 为主，还原性远没有甲烷—氨混合气那么强；而在这种“温和”大气里放电，氨基酸的产量要低得多。米勒实验因此被批评为“模拟了一个并不存在的地球”。

故事还没完。米勒去世后，他的学生杰弗里·巴达在实验室里翻出了 1950 年代封存的原始样品瓶——当年做完实验、却从未分析过的那种。2008 年，巴达用现代仪器重新分析，发现其中一个模拟“火山喷发”（气体里加了 \chem{H_2S}）的实验，竟生成了 22 种氨基酸，比当年发表的还多。这提示：即使全球大气不合适，火山附近、热泉口这类\textbf{局部环境}依然可能是高效的“零件工厂”。

这个故事值得记住三点：好实验也可能建立在后来被修正的假设上；假设被修正，不等于实验白做——它证明了“无机物变有机零件”这条路本身走得通；而科学的诚实体现在，他们连五十年前的旧样品都肯翻出来重新查一遍。
\end{zhengju}

\section{规则层：一锅稀汤为什么还不是生命}

零件撒进海洋，得到的只是一锅被无限稀释的“有机汤”。从这锅汤到“一个细胞”，中间隔着三道难关，每一道都得靠你已经学过的规则去翻。

\textbf{难关一：浓度。}两个分子要反应，先得撞上。可海洋太大了：把零件撒进去，就像往西湖里倒一勺糖，谁也碰不见谁。稀到什么程度？算一笔账——

\begin{qiaoliang}
想象一个直径 1 微米的小泡（和今天的细菌差不多大），体积约 $5\times 10^{-16}$ 升。假如里面某种零件的浓度是每升 1 毫摩尔（$10^{-3}$ 摩尔，已经不算特别稀），这个小泡能装下几个这种分子？

物质的量：$5\times 10^{-16}\times 10^{-3}=5\times 10^{-19}$ 摩尔；乘以每摩尔 $6.02\times 10^{23}$ 个，得到约 $3\times 10^{5}$ 个——三十万个。听起来不少？但一条能自我复制的分子链，可能需要几百个零件按正确顺序首尾相连；而没有边界的大海“稀释汤”里，真实浓度可能比这还要低成千上万倍。没有浓缩机制，撞出复杂分子的概率低到令人绝望。
\end{qiaoliang}

好在早期地球并不缺浓缩的办法，而且都是你在生活中见过的：潮汐池被太阳晒干，溶质越来越浓——和你把一碟盐水放在窗台上、最后结出盐晶是同一回事（第 25 章）；海水结冰时，杂质被挤进冰缝里残留的盐水中，浓度飙升——冻过一瓶糖水再化开，头一口准是齁甜的；热泉矿物里那些微米级的孔隙，则像无数口小锅，各煮各的汤。

\textbf{难关二：边界。}第 45 章说过，生命系统的第一个特征是“有里面和外面”。没有边界，任何有用的分子刚一合成，就被大海冲散了。好消息是：边界可以不请自来。第 47 章讲过，磷脂和更简单的脂肪酸都是“一头亲水、一头躲水”的分子，撒进水里会自动聚成双层膜的小球——小泡。这是疏水效应（第 22、23 章）的纯物理结果，不需要任何生命帮忙。在今天的实验室里，杰克·绍斯塔克等人的团队发现：脂肪酸小泡不但能自己形成，喂给它新的脂肪酸分子还会长大，受到水流剪切时会一分为二——一座会“生长、分裂”的膜房子，整个过程没有任何基因参与。

\begin{center}
\begin{tikzpicture}
\draw[thick] (0,0) circle (1.5);
\draw[thick] (0,0) circle (1.25);
\foreach \a in {0,30,60,90,120,150,180,210,240,270,300,330} {
  \fill (\a:1.375) circle (0.06);
}
\draw (0.15,-0.55) -- (0.45,-0.35) -- (0.2,-0.1) -- (0.5,0.15);
\fill (0.15,-0.55) circle (0.05);
\fill (0.45,-0.35) circle (0.05);
\fill (0.2,-0.1) circle (0.05);
\fill (0.5,0.15) circle (0.05);
\node at (0,0.6) {\small 内部的水溶液};
\draw[->,thick] (2.6,0) -- (1.65,0);
\node[right] at (2.65,0) {\small 脂肪酸分子};
\end{tikzpicture}
\end{center}

上面这张图里，两层圆圈代表自动围成双层的脂肪酸膜，折线代表被包进内部的核酸链——纯属人为的表示法，分子既不真是圆点，也不真是线段。

\begin{shiyan}
\textbf{材料}：透明杯子两个、水、食用油、洗洁精、筷子。

\textbf{步骤}：① 两个杯子各装半杯水，各滴入 5 滴食用油，用筷子各搅 10 秒，静置观察。② 往第二个杯子里加一滴洗洁精，再搅 10 秒，静置对比。

\textbf{观察点}：没加洗洁精的杯子里，油很快聚成大大小小的圆球，浮上水面——油分子“躲水”，自动抱团，缩小和水的接触面。加了洗洁精的杯子里，油被打散成更多、更小、更稳定的微小油滴，久久不肯合并——洗洁精分子一头亲水、一头亲油，围在油滴外面，相当于一层“分子膜”。这正是“两亲分子自动围出边界”的缩影：小泡的形成不靠图纸，只靠分子自身的性格和水分子的排斥规律。

\textbf{安全提示}：材料全是厨房级别的，无危险；做完的液体倒入水槽冲走即可。
\end{shiyan}

\textbf{难关三：复制。}一栋膜房子再漂亮，不能复制，就只是一次性的烟火。要走向生命，汤里必须出现某种\textbf{能携带信息、又能照着信息复制自己}的分子。而且复制还必须有误差——这一点听起来反直觉，却至关重要：完美复制只有重复，没有变化；带着误差的复制才产生变异，而变异是自然选择的原料（第三册会详细展开）。抄写员抄书难免抄错字，复印件复印次数多了会糊——分子复制，同理。

\section{符号层：DNA 出场前，先认识它的堂兄}

今天你的细胞里，信息存在 DNA 里，干活的主要是蛋白质，RNA 在中间当信使和工具（第三册会细讲这条流水线）。但仔细看看 RNA 这个“中间人”，你会发现它有点不一般。

RNA 和 DNA 一样，是一条核苷酸长链：每个核苷酸由磷酸、核糖和四种碱基之一组成，四种碱基记作 A、U、G、C。链与链之间遵守严格的配对规则：A 只认 U，G 只认 C——靠的是氢键的形状匹配（第 22 章），像两排凹凸互补的拉链齿。正是这条配对规则，让一条链能当“模板”，指导合成出与它互补的另一条链。这就是复制的化学本质；第三册第 66 章会把 DNA 的这套结构彻底拆开看。

把核苷酸连成链，每接上一节就要脱去一个水分子——和氨基酸连成蛋白质时形成肽键（第 48 章）一样，都是缩合反应，都需要能量推一把。这就是难关三的收费处：模板配对是“信息”问题，脱水成链是“能量”问题，两边都得解决。

而 RNA 真正让人坐直身子的事实在于：\textbf{它不只是信息载体，它还会干活。}这就是下一节的主角。

\section{RNA 世界：一个漂亮但未被证实的假说}

1980 年代，两个独立的发现震惊了生物学界。托马斯·切赫在研究四膜虫时发现，一段 RNA 能在没有任何蛋白质帮忙的情况下，把自己中间的一段“剪”掉、再把两头接好；西德尼·奥尔特曼则证明，一种叫 RNase P 的酶里，真正干催化活的是它的 RNA 部分，蛋白质只是配角。会催化的 RNA 被命名为\textbf{核酶}，两人因此获得 1989 年诺贝尔化学奖。更惊人的还在后面：2000 年前后，核糖体的高分辨率结构显示，细胞里制造蛋白质的那台大机器，其“焊接”氨基酸的核心部位也是由 RNA 构成的——你体内每一条蛋白质，归根结底都是 RNA 催化出来的。

这些发现催生了一个大胆的想法：既然 RNA 既能像 DNA 一样存信息，又能像酶一样催化反应，那么\emph{在 DNA 和蛋白质分工之前，会不会存在一个“RNA 世界”}——第一批能自我复制的系统，从复制到催化，全靠 RNA 一个分子扛？1986 年，沃尔特·吉尔伯特为这个假说命了名。按照这幅图景，DNA 稳定、适合长期存档，蛋白质能干、适合当机器，但它们都是后来才接管事务的“专业部门”；创业初期，是 RNA 一人分饰两角。

现在，请把下面这句话读两遍：\textbf{RNA 世界是一个假说，不是一段已经被重演的历史录像。}它有真实证据支撑——核酶确实存在，核糖体的核心确实是 RNA——而且能统一解释很多事实，是目前最有影响力的图景之一。但实验室至今只重演过它的\textbf{片段}：化学家在模拟早期条件下合成出了核苷酸的部分原料，让无酶的 RNA 复制跑过短短几轮，造出过会长大分裂的小泡。把这些片段拼成“从无机物到活细胞”的完整链条？还没有人做到。RNA 的零件在早期地球上怎样批量制造、没有酶的复制怎样做到既够长又够准、第一批 RNA 究竟从哪儿来——每一环都还没有闭合。

\section{从化学演化到达尔文演化：门槛在哪里}

退一步看全局。从一锅有机汤到第一批细胞，科学家对“需要凑齐哪几件事”意见大致一致——但对“具体怎样凑齐、按什么顺序凑齐”，还在激烈争论。先看看这四块门槛：

\begin{enumerate}[itemsep=2pt]
\item \textbf{个体化的边界}：膜小泡把“一个系统”从环境里圈出来，让有用的分子能积累、能“私有”，而不是随海浪漂走。
\item \textbf{从环境获取能量和材料的通路}：也就是最原始的“代谢”——哪怕只是热泉口氢气和 \chem{CO_2} 的反应，顺路推着别的反应走（第 50 章讲过的放能—吸能偶联的雏形）。
\item \textbf{能复制、且复制有误差的信息载体}：比如一段 RNA。
\item \textbf{误差影响复制的成功率}：有的小泡里，分子复制得快一点，或者膜结实一点、漏得慢一点。
\end{enumerate}

第四件事一到位，自然选择就\textbf{自动}启动了——不需要谁设计，也不需要谁指挥。这是纯粹的统计规律（第 28 章）：复制成功率高的版本，后代占的比例必然越来越大，一代又一代，差距被无情放大。走到这一步，“化学演化”（分子在能量和概率规则下越变越复杂）就跨过了门槛，变成“达尔文演化”（变异被筛选、被积累）。这颗星球从此有了“历史”——此后近 40 亿年的故事都是这条逻辑的展开，第三册会一路讲到进化本身。

但这里有一个至今悬而未决的大问题，学界分成好几派：\textbf{代谢和复制，谁先谁后？}“复制先行”派（大致就是 RNA 世界路线）认为，先有了能复制的分子，小泡和代谢是后来配上的；“代谢先行”派认为，热泉口的化学能先驱动起一张简单的反应网络——类似第 52、53 章那些代谢路径的“原始版”——遗传分子是这张网络里后来“结晶”出来的产物。双方手里都有实验证据，也都有解释不了的缺口；还有人认为，答案可能不是二选一，而是两者纠缠着共同起步。诚实的回答是：\textbf{现在还不知道。}这不是科学的失败，这正是科学正在工作的样子。

\section{边界层：这是一张拼图，不是一盘录像带}

本章的所有图景——米勒的天空、热泉的孔隙、自组装的小泡、RNA 世界——正确的理解方式是：它们是“\textbf{与现有证据相容的可能路径}”，而不是“当年就是这样发生的”。理由有三。

第一，直接证据已被地球销毁，我们永远无法回放实况。第二，每条假说都有没补上的洞：RNA 世界缺“无酶条件下可靠复制”这一环；热泉路线缺“反应网络怎样变成遗传系统”的桥梁；甚至有人提出，RNA 之前可能还有更简单的遗传分子（所谓“前 RNA 世界”）——连主角是谁都还没定下来。第三，这个领域的图景会被新证据反复改写：米勒的大气模型被修正，就是前车之鉴。

所以本章模型“好用”的地方在于：它告诉你生命起源\textbf{不违反}任何化学和物理规律，而且每一步——自组装、催化、模板复制、自然选择——单拎出来都能在实验室里看到。“不好用”的地方在于：你不能拿它当历史课本来背。如果哪天你在某本书里读到“科学家已经破解了生命起源之谜”，请想起这一章，笑一笑，翻过去。

\begin{wuqu}
“让原子随机碰撞组装出一个细胞，就像龙卷风卷过垃圾场、当场装出一架波音 747——概率上不可能，所以生命起源必有神秘力量插手。”

这个比喻错在两点。第一，没有人主张细胞是一步撞出来的。真实的图景是无数\textbf{小台阶}：小泡自组装、核苷酸逐步形成、短链先开始复制、复制体再被筛选——每一步都是普通化学，概率并不低；而且每一步的产物会被化学规律和自然选择“留下”，不是撞完就散。龙卷风的比喻，把“有记忆、有筛选的积累过程”偷换成了“一步到位的抽奖”。第二，它暗示有序结构的出现违反热力学第二定律——并没有。地球是开放系统，太阳和地球内部源源不断输入能量；局部变得更有序，只要整体的熵在增加就合法。雪花结晶、盐粒析出、婴儿发育，全是局部有序化的日常例子（第 28 章）。747 的比喻拿来提醒“别想一步到位”还算形象，拿来否定整个研究领域，就是用错了地方。
\end{wuqu}

\section{回到开场：鹅卵石和虫子之间隔着什么}

现在可以回答开场的问题了。鹅卵石和小虫，确实出自同一盒积木——而且小虫身体里的每个碳原子、每个铁原子，都和石头里的一样古老（第 16 章：元素在恒星里造好，地球只是继承）。区别不在积木，在\textbf{组织方式和历史}：虫子体内有膜围出的亿万个“里面”，有能量流日夜冲刷、维持着这些“里面”不塌（第 50、53 章），还有一套抄写了近 40 亿年、带着抄错痕迹的信息档案。石头呢？一条都不占：没有“里面”和“外面”的分别，没有能量流过它，更没有任何东西在它身上复制自己。

所以“活”不是某一种原料，而是一种\textbf{运转方式}。这种运转方式最初怎样从石头和海水里启动——本章给你的，是目前能查到的最好的候选答案，以及最诚实的“还不知道”。

\section{继续深挖：一个“还不知道”养活了多少科学}

你可能没想到，这个连最终答案都没有的领域，已经结出了实实在在的果子。医生手里的许多药物和疫苗，靠\textbf{脂质体}送货：用人工磷脂小泡把药物分子包起来，让它安全抵达目标细胞；mRNA 疫苗的核心技术，就是把一段 mRNA 装进脂质纳米颗粒。看出来了吗——“脂质小泡包着核酸”，这不就是本章那张原始细胞草图的现代技术版吗？研究起源的科学家造“原始小泡”，药剂师造“送货小泡”，用的是同一条自组装原理。

合成生物学则在从另一头逼近这个问题：科学家造出了基因组被精简到极限的人造细胞，数清了“一个活系统最少需要多少个零件”；也有团队让 RNA 分子群体在试管里一代代进化——复制、变异、选择，全部真实发生。这些都不是“重演起源”，而是在测绘“可能性空间的地图”。

这套图景甚至改写了人类寻找外星生命的方式：如果生命真能在热泉边靠化学能起步，那么火星的古老湖床、木卫二冰壳之下的海洋，都值得去翻一翻。

至于第一个真正的细胞长什么样、它怎样一步步变成今天你体内这座有细胞核、有线粒体、有骨架的繁华城市——下一章，我们就钻进现代细胞内部去看看。你会在那里找到远古留下的活化石：每一台核糖体的核心，至今仍是 RNA。

\begin{tiaozhan}
还有一个被本章悄悄跳过的谜：手性。第 42 章讲过，许多分子像左右手一样互为镜像，而地球生命极度“偏食”——蛋白质几乎只用左旋氨基酸，核酸只用右旋糖。可是默奇森陨石里的氨基酸，左旋右旋几乎各占一半，说明化学本身并不挑边。那么生命为什么、又是怎样“选边”的？有人认为纯属偶然：第一个成功的复制系统恰好用了一边，后代就被“锁死”在这条祖传流水线上——就像铁路轨距，一旦铺定，改起来代价太大；也有人找到了能放大微小不对称的物理化学机制（某些晶体的表面、星际空间的圆偏振光），让“选边”可能并非纯粹掷硬币。目前仍无定论。把它当成一个提醒吧：连“生命的分子为什么长这只手”都还没有答案，起源研究离收官还远得很。
\end{tiaozhan}

\section*{练一练}

\begin{sikaoti}
\item 画一张“从零件到原始细胞”的物质流图：起点是 \chem{CO_2}、水、\chem{N_2} 等无机小分子，终点是“膜小泡里装着可复制的分子”。在每条箭头上注明推动这一步的能量来源（光、热、化学能）或规则（疏水效应、碱基配对）。
\item 假如把米勒实验的大气换成 \chem{CO_2}、\chem{N_2} 和水蒸气，氨基酸产量明显下降。用第 27、29 章的语言解释：改变反应物，为什么会同时影响“能生成多少”和“生成得多快”？这两件事分别由什么决定？
\item 有人说：“晒干的潮汐池把分子从稀汤浓缩成浓汤，这是熵减小，违反热力学第二定律。”用第 28 章的知识，写三句话反驳他。
\item 绍斯塔克的脂肪酸小泡会“长大、分裂”，但它还不算生命。对照本章列的四块门槛，指出它还缺哪几块，并说明为什么缺了它们，自然选择就无法启动。
\item “RNA 世界是一个假说”——这里的“假说”和日常说的“我有个假说”（约等于瞎猜）有什么区别？从本章的证据故事里举两个例子，说说科学假说和瞎猜之间的界线划在哪里。
\item 如果在木卫二的冰下海洋里发现了热泉和简单有机物，但没有发现任何会复制的分子——按照本章的图景，这说明那里走到了四块门槛中的哪一步？
\end{sikaoti}
